\documentclass{article}
\usepackage{pathfinder-readable}

\title{Quantile Advice and Distributed CVaR Optimisation: A Paused Connexion}

\begin{document}
\maketitle

\begin{abstract}
This note considers a survey of learning-augmented algorithms, Q, alongside a paper on distributed
optimisation of conditional value at risk, P. The thread asked whether a predicted tail threshold could
reduce the loss queries used by P. A threshold prediction alone does not reduce P's batch size, but the
peers derived a different, one-query method and a bound for one agent. They did not prove that advice
improves that method or that its bound extends to P's network. The verifier paused the thread pending a
causal, cost-aware comparison with a method that receives no advice.
\end{abstract}

\section{Introduction}

Q surveys algorithms that use fallible predictions while retaining guarantees when predictions fail
\cite{Q}. It asks what is predicted, how error is measured, and what it costs to obtain and use a
prediction. Its consistency and robustness language is framed mainly through cost ratios. Those questions
guide the present enquiry, but its ratio guarantees do not directly describe a risk-minimisation gap.

P studies agents that jointly minimise the average of their local conditional values at risk (CVaR) over
a changing communication network \cite{P}. CVaR measures the mean of the adverse tail of a loss
distribution. Each agent sees noisy loss values, rather than a risk gradient. P evaluates a batch of
losses at a perturbed decision, forms an empirical CVaR, and uses that estimate in a distributed update.
It proves consensus and risk bounds under its stated convexity and regularity assumptions.

The proposed connexion was to give each agent a predicted value at risk (VaR), a threshold near the start
of the adverse tail, and spend fewer loss queries. The peers first tested that idea against P's actual
sampling procedure. When it failed, they examined whether keeping the threshold as part of the
optimisation state could make advice useful in a different method.

\section{What was done}

The peers first checked where P spends its queries. P obtains independent full-loss observations before
minimising its empirical CVaR over a scalar threshold. A predicted threshold might speed that scalar
minimisation, but does not supply missing tail observations or remove the batch's sampling error. The
peers therefore stopped treating a VaR seed as a query-saving change to P's algorithm.

They next used the threshold representation of CVaR that P already states. For an agent \(i\), let \(x\)
be its decision, \(J^i(x,\xi^i)\) its random loss, \(\alpha_i\) its upper-tail probability, and \(t_i\) a
threshold. With \((a)_+=\max\{a,0\}\), the lifted objective is
\[
F_i(x,t_i)=t_i+\frac{1}{\alpha_i}
\mathbb E\bigl[(J^i(x,\xi^i)-t_i)_+\bigr].
\]
Minimising \(F_i\) over \(t_i\) gives the agent's CVaR. Instead of recomputing empirical CVaR from a
batch, the proposed method keeps \(t_i\) as a local state and updates it with \(x\). It mixes decisions
between neighbours as P does; the network update itself was not proved.

For each proposed update, the agent drew one fresh loss at a randomly perturbed decision. The peers used
the observed excess over \(t_i\) for the decision update and an indicator of whether the loss exceeded
\(t_i\) for the threshold update. Subtracting the known threshold from the decision estimator removed a
direction-independent term. It did not change the estimator's mean, because the perturbation direction
has mean zero. This construction used one loss evaluation per agent per step and avoided the error from
plugging a finite-batch CVaR estimate into the gradient. The detailed estimator and its
conditional-expectation check are in \texttt{ada/quantile\_lift\_note.md}.

The peers then bounded the estimator's second moment, proved a projected-update bound for one agent, and
corrected its smoothing term. They compared sufficient query bounds under homogeneous tuning. Finally,
they tested whether the variance improvement could actually be credited to predictions, and whether the
one-agent argument carried over to P's network. Both questions remained open.

\section{Results supported by the checks}

The one-query construction has a prediction-independent variance bound. To state the sharper conditional
form, let \(d\) be the dimension of \(x\), \(\delta_i\) the perturbation radius, and \(U_i\) a bound on
the absolute loss and on the clipped threshold. Let \(p_{i,k}\) be the probability, conditional on the
past, that the fresh perturbed loss at step \(k\) exceeds the current threshold. The peers showed that
the decision estimator's conditional second moment is at most
\[
\frac{4d^2U_i^2p_{i,k}}{\alpha_i^2\delta_i^2}.
\]
The threshold estimator's conditional second moment equals \(1-2p_{i,k}/\alpha_i+p_{i,k}/\alpha_i^2\).
Since \(p_{i,k}\leq 1\), both expressions have a bound even for arbitrary advice. If thresholds remain
calibrated so that \(p_{i,k}\) is of order \(\alpha_i\), their risk-level dependence improves from order
\(\alpha_i^{-2}\) to order \(\alpha_i^{-1}\), with other quantities fixed. This is a condition on the
whole trajectory, not a proved effect of a predictor. The derivation appears in the peer check
\texttt{ada/quantile\_lift\_note.md}.

For one agent, the peers applied convexity and the projected stochastic-subgradient recursion to the
joint decision and threshold. They bounded the expected CVaR gap of the step-size-weighted average by
three terms: the initial squared distance from a shrunken comparator and its optimal smoothed threshold,
divided by the sum of step sizes; a step-size-weighted sum of the two second moments above; and a
smoothing and domain-shrink term. Here the comparator is an optimal decision moved into the shrunken
feasible domain required for perturbation. If \(L\) bounds the samplewise loss change per unit decision
distance, \(\delta\) is the smoothing radius, \(D_x\) bounds the comparator's distance from the origin,
and \(r\) is P's interior-ball radius, the last term is at most \(L\delta(2+D_x/r)\). The peers obtained
this sharper term by choosing the optimal threshold \emph{after} smoothing at the shrunken comparator.
The complete inequality and assumptions are in \texttt{ada/quantile\_lift\_note.md}. In particular, this
is a one-agent additive gap bound, not a distributed theorem or an advice advantage.

With homogeneous tuning, the one-agent bound gives a sufficient query count of order
\(\alpha^{-2}\varepsilon^{-4}\) for target gap \(\varepsilon\) in the worst case, or
\(\alpha^{-1}\varepsilon^{-4}\) if calibrated tail occupancy persists. Here \(\alpha\) is the common tail
probability. For context, tuning P's displayed polynomial-step guarantee gives \(\widetilde
O(\alpha^{-2}\varepsilon^{-10})\) loss queries; reoptimising its underlying ergodic inequality gives the
different comparison \(\widetilde O(\alpha^{-2}\varepsilon^{-8})\). The tilde hides logarithmic factors.
These are sufficient upper-bound calculations under different analyses, not lower bounds or a matched
comparison of algorithms. The correction to P's displayed rate is recorded in the peers' checks and the
ledger.

\section{Objections and limits}

Low tail-activation probability does not establish that advice helped. The prediction-free threshold
\(t_i=U_i\) starts with \(p_{i,0}=0\). A threshold set too high can also make the decision estimator
quiet while worsening the lifted objective. A meaningful advice error must therefore reflect quantile
accuracy or excess lifted objective, and any gain must be measured against a matched rule that receives
no prediction and pays the same query budget, including predictor calls.

An accurate VaR at the observable initial decision need not equal the optimal smoothed threshold at the
unknown comparator decision. It therefore need not remove the initial-distance term in the one-agent
bound. Predictions issued at later decisions might help, but resetting thresholds from them would require
a separate tracking argument. The peers did not supply one.

The network step poses a further obstacle. With a threshold held fixed, the lifted local objective can be
\(L_i/\alpha_i\)-Lipschitz in the decision, where \(L_i\) is the agent's samplewise loss Lipschitz
constant. CVaR after minimising over thresholds has the sharper \(L_i\) bound. P's disagreement proof
uses local-objective regularity, and the agents' thresholds vary, so the one-agent rate cannot simply be
copied into a distributed claim. The full network recursion, matched-budget tests, and a broader
prior-work check remain outstanding.

The peers also checked two cited precedents. Soma and Yoshida already use joint decision-threshold CVaR
optimisation with one sample in a first-order setting \cite{somaYoshida}. Wang, Shen, and Zavlanos study
residual feedback for one-point CVaR estimates \cite{wangShenZavlanos}. Those checks limit novelty
claims; the peers' searches did not establish that an identical distributed result is absent.

\section{Why the thread stopped}

The verifier paused the thread because the supported one-query argument did not establish a benefit from
quantile advice, and the distributed guarantee was unproved. The smallest step that could restart the
connexion is a one-agent theorem specifying when a prediction is available, what it costs, and how its
error affects an additive CVaR gap against a matched method with no advice. If that comparison shows a
benefit, the network extension becomes a concrete next question.

\bibliographystyle{plainurl}
\bibliography{references}
\end{document}
